\chapter{Contexte général et étude de l’existant}

\minitoc

\newpage

\section{Introduction}

Ce chapitre a pour objectif de présenter le cadre général du projet de fin d’études. 
Nous commençons par une présentation de l’entreprise d’accueil ainsi que de ses principaux domaines d’activité. 
Ensuite, nous étudions les solutions existantes dans le domaine concerné afin d’identifier leurs limites et les améliorations possibles. 
Enfin, nous présentons les exigences du projet ainsi que la méthodologie de travail adoptée durant sa réalisation.

% =====================================================




\section{Présentation de la société}

Cette section est consacrée à la présentation de la société d'accueil au sein de laquelle ce projet a été réalisé. Nous y décrivons son historique, son positionnement dans l'écosystème technologique tunisien ainsi que ses principaux domaines d'intervention.

% =====================================================

\subsection{Présentation générale}

Standard Sharing Software, désignée sous l'acronyme \textbf{3S} \cite{site3s}, est une entreprise tunisienne fondée en 1988, spécialisée dans l'intégration de solutions technologiques et les services informatiques à forte valeur ajoutée. Forte de plus de 35 ans d'expérience, elle s'est imposée comme un acteur de référence sur le marché tunisien et africain, en accompagnant les organisations publiques et privées dans leurs projets de modernisation numérique.

3S se distingue par une approche globale qui couvre l'ensemble du cycle de vie des projets IT : de l'étude et de l'audit des besoins jusqu'au déploiement et au pilotage des solutions mises en place. Cette capacité à intervenir de bout en bout, combinée à des partenariats avec des éditeurs et constructeurs de renommée mondiale, lui permet de proposer des réponses adaptées aux exigences les plus complexes.

\subsection{Positionnement au sein du Groupe 3S}

3S est la filiale fondatrice du \textbf{Groupe 3S} \cite{site3sgroup}, un groupe technologique tunisien structuré autour de deux pôles complémentaires. Le \textbf{pôle IT} regroupe les filiales dédiées aux technologies de l'information : \textbf{3S} (intégration IT), \textbf{GlobalNet} (accès Internet grand public), \textbf{GlobalNet Business} (services cloud et sécurité pour entreprises), \textbf{Advancia Training} (formation et certification IT), \textbf{Smartup} (édition de progiciels), \textbf{Hexabyte} (connectivité et solutions high-tech), \textbf{Cyberesa} (solutions e-tourisme) et \textbf{Pro2C} (centre de contact multicanal). Le \textbf{pôle Cleantech} quant à lui rassemble des entités orientées développement durable, telles que \textbf{Africa Solar} (énergies renouvelables), \textbf{CandyLed} (éclairage LED), \textbf{BioPack} (bioplastiques) et \textbf{Océane} (construction navale).

% =====================================================



\subsection{Domaines d'activité de 3S}

En tant que pôle IT du groupe, 3S intervient sur l'ensemble de la chaîne de valeur liée à la transformation numérique des entreprises. Ses domaines d'activité s'articulent autour de six axes principaux :

\begin{itemize}

    \item \textbf{Intégration des infrastructures IT :}
    3S accompagne ses clients dans la conception et le déploiement d'environnements informatiques robustes et évolutifs, conformes aux standards internationaux, depuis l'audit initial jusqu'à la mise en production ;

    \item \textbf{Réseaux et connectivité :}
    l'entreprise propose des solutions d'interconnexion et de télécommunication garantissant des échanges fiables et sécurisés entre les différents composants des systèmes d'information ;

    \item \textbf{Datacenter et continuité de service :}
    virtualisation, sauvegarde, stockage et haute disponibilité constituent les piliers de l'offre datacenter de 3S, visant à assurer la résilience et la protection des données critiques ;

    \item \textbf{Cloud computing :}
    la société déploie des architectures cloud privées, hybrides et publiques, ainsi que des services d'hébergement managés, répondant aux besoins croissants des entreprises en matière de flexibilité et d'agilité ;

    \item \textbf{Cybersécurité :}
    3S propose une offre complète de sécurité informatique intégrant des pare-feux nouvelle génération, des systèmes SIEM, une protection Anti-DDoS, des solutions de gestion des identités et des accès (IAM), ainsi que des centres opérationnels SOC et NOC dédiés à la surveillance et à la réponse aux incidents ;

    \item \textbf{Communication unifiée et collaboration :}
    l'entreprise intègre des outils de communication unifiée permettant d'améliorer la coordination et la productivité des équipes, en facilitant les échanges au sein des organisations.

\end{itemize}

Cette palette de compétences, combinée à l'écosystème du Groupe 3S, permet à 3S de se positionner comme un partenaire technologique global, capable de répondre aux enjeux actuels de la numérisation, de la sécurisation des systèmes d'information et de l'évolution des infrastructures modernes.

% =====================================================





\section{Contexte du projet}

Ce projet intitulé \textit{« Conception et développement d'une plateforme intelligente de montée 
en compétences et de mentorat en entreprise, pilotée par l'intelligence artificielle »} a été 
réalisé au sein de la société Standard Sharing Software (3S) dans le cadre du Projet de Fin 
d'Études en vue de l'obtention du diplôme national d'ingénieur.

% =====================================================

\subsection{Problématique}

Dans un environnement professionnel en constante évolution, les entreprises peinent à assurer 
une montée en compétences efficace de leurs collaborateurs, qu'il s'agisse d'employés en poste 
ou de stagiaires en phase d'intégration. Les dispositifs de formation classiques montrent 
aujourd'hui leurs limites : les études indiquent que seulement 10 à 15\% des cours en ligne 
sont effectivement complétés, révélant un problème profond de motivation et d'engagement des 
apprenants face à des contenus statiques et non personnalisés.

À cela s'ajoute l'absence quasi totale de mécanismes de coaching et de mentorat structurés 
au sein des plateformes existantes. Les stagiaires et les nouveaux collaborateurs se retrouvent 
souvent livrés à eux-mêmes, sans suivi individualisé ni accompagnement adapté à leur profil 
et à leur progression. Les outils disponibles sur le marché restent cloisonnés, rigides et peu 
intelligents : ils ne permettent pas d'identifier les lacunes en temps réel, de recommander des 
contenus pertinents ni de mettre en relation efficacement un apprenant avec un mentor qualifié.

La question centrale de ce projet est donc la suivante : \textit{comment concevoir une 
plateforme capable de rendre la formation en entreprise plus engageante, mieux suivie et 
intelligemment accompagnée, notamment pour les employés et stagiaires en phase de montée 
en compétences ?}

% =====================================================

\subsection{Objectif}

Ce projet vise à concevoir et développer une plateforme web intelligente qui centralise la 
gestion des formations, le suivi des compétences et le mentorat au sein de l'entreprise. 
L'objectif est de proposer à chaque apprenant un parcours personnalisé, soutenu par des 
recommandations intelligentes et un coaching adapté à son profil et à sa progression.

La plateforme intègre également des mécanismes de gamification pour renforcer la motivation 
des utilisateurs et améliorer le taux de complétion des formations. Sur le plan architectural, 
elle est conçue selon une approche multi-tenant et modulaire, permettant à plusieurs entités 
du Groupe 3S de l'adopter de manière indépendante au sein d'une infrastructure commune.

% =====================================================

\subsection{Solution proposée}

Pour répondre à ces problématiques, nous proposons une plateforme web intelligente intégrant 
la gestion des utilisateurs, des formations, des compétences et des parcours d'apprentissage, 
ainsi qu'un module de mentorat structurant la relation entre mentors et apprenants.

L'intelligence artificielle est au cœur de la solution : des agents IA dotés d'une mémoire 
contextuelle assurent des recommandations personnalisées, une assistance continue et un 
coaching adapté à chaque utilisateur. Des fonctionnalités de gamification complètent cette 
expérience en encourageant l'engagement et la progression des apprenants tout au long de 
leur parcours.

L'architecture retenue est de type multi-tenant, permettant aux différentes filiales du 
Groupe 3S de disposer chacune d'un espace indépendant au sein d'une infrastructure partagée, 
garantissant ainsi l'isolation des données, la flexibilité de configuration et la scalabilité 
de la solution.




\section{Étude et analyse de l'existant}


Avant d'entamer la conception de notre solution, nous avons mené une étude approfondie 
des plateformes existantes dans le domaine de la formation et du développement professionnel. 
Cette analyse se concentre sur les plateformes qui intègrent des approches modernes telles 
que le coaching intelligent piloté par l'IA, la gamification, la génération automatique de 
cours et l'utilisation d'agents conversationnels. Elle s'intéresse également à des critères 
pratiques tels que la confidentialité des données et le modèle tarifaire de chaque solution. 
L'objectif est d'évaluer dans quelle mesure ces plateformes parviennent à améliorer la 
motivation des apprenants, à renforcer leur engagement et à leur offrir un accompagnement 
personnalisé et efficace. Cette étude s'articule autour de deux catégories : les solutions 
commerciales disponibles sur le marché, et les solutions non commerciales issues de la 
communauté open source ou de la recherche académique.


% =====================================================

\subsection{Solutions commercialisées}

Plusieurs plateformes commerciales ont intégré des fonctionnalités avancées d'intelligence 
artificielle pour améliorer l'expérience d'apprentissage en entreprise. Nous présentons 
ci-dessous les solutions les plus représentatives du marché actuel, en nous concentrant 
sur leurs approches pédagogiques, leurs fonctionnalités clés et leur positionnement.



% =====================================================
\subsubsection{Google Skills}

Lancée en octobre 2025, Google Skills \cite{googleSkillsFeatures} est une plateforme d'apprentissage 
unifiée regroupant près de 3\,000 cours, labs et certifications issus de l'écosystème Google 
(Google Cloud, DeepMind, Grow with Google et Google for Education). Elle cible des profils 
allant du débutant au développeur confirmé, avec un accent fort sur l'IA générative et les 
technologies cloud. La plateforme mise sur la gamification pour stimuler l'engagement des 
apprenants à travers des ligues de points, des badges de compétences, des streaks 
d'apprentissage et des tableaux de classement par équipe. Pour les entreprises clientes de 
Google Cloud, un accès organisationnel gratuit est disponible via Google Skills for Orgs, 
permettant d'assigner des parcours aux équipes et de suivre leur progression via des 
tableaux de bord analytiques.

% =====================================================



\subsubsection{Duolingo}
Duolingo \cite{duolingo} est une plateforme d’apprentissage des langues largement reconnue pour son approche basée sur la gamification. Elle intègre des mécanismes tels que les points d’expérience, les séries de connexions quotidiennes (streaks) et les classements, afin de renforcer l’engagement et la régularité des apprenants. Son modèle pédagogique repose sur des exercices courts et interactifs, favorisant un apprentissage progressif et motivant.

Ces dernières années, la plateforme a renforcé son orientation vers l’intelligence artificielle en automatisant une grande partie de la création de ses contenus pédagogiques, avec le déploiement de nombreux cours générés par IA. Par ailleurs, son offre premium Duolingo Max intègre des fonctionnalités avancées basées sur des modèles de langage, notamment des agents conversationnels permettant de simuler des dialogues réalistes, ainsi qu’un moteur adaptatif qui ajuste dynamiquement la difficulté des exercices en fonction du niveau et des performances de l’apprenant.



\subsubsection{360Learning}
360Learning \cite{learning360} est une plateforme d’apprentissage collaborative destinée aux entreprises, basée sur un modèle où les employés deviennent eux-mêmes créateurs de contenus pédagogiques. Elle favorise ainsi une approche d’apprentissage pair-à-pair et une production de formations directement issues des experts métiers.

La plateforme se distingue par son fort niveau d’engagement des apprenants, lié à son intégration dans le travail quotidien et à son approche collaborative. Elle intègre également des fonctionnalités d’intelligence artificielle, notamment des outils de recommandation de contenus et de création assistée permettant de transformer rapidement des documents (PDF, PowerPoint, Word) en modules de formation structurés.

\subsubsection{CYPHER Learning}CYPHER Learning \cite{cypherLearning} est une plateforme LMS et LXP alimentée par l'IA, 
positionnée comme solution tout-en-un pour la création, la gestion et la personnalisation 
des parcours d'apprentissage. Son outil phare, CYPHER Copilot, permet de générer 
automatiquement des cours complets, des quiz et des évaluations à partir de ressources 
fournies par l'organisation, avec un mécanisme de vérification de l'exactitude des 
contenus générés (AI Crosscheck). La plateforme propose également une gamification pilotée 
par l'IA, adaptant les points, badges et niveaux au profil de chaque apprenant, ainsi 
que des parcours adaptatifs qui s'ajustent dynamiquement selon les résultats obtenus.

% =====================================================





\subsection{Solutions non commercialisées}
Des initiatives issues de la communauté open source et de la recherche académique proposent des alternatives libres, transparentes et hautement personnalisables aux plateformes propriétaires. Ces solutions présentent un intérêt particulier pour les organisations qui souhaitent conserver la maîtrise de leurs données ou expérimenter des approches pédagogiques innovantes.

% =====================================================
\subsubsection{OpenMAIC}

OpenMAIC \cite{openmaic} (Open Multi-Agent Interactive Classroom) est une plateforme open source développée par l’équipe de recherche MAIC de l’Université de Tsinghua et publiée en 2025. Elle propose une approche pédagogique basée sur des systèmes multi-agents, allant au-delà des chatbots éducatifs classiques.

À partir d’un document ou d’un sujet décrit en langage naturel, la plateforme génère automatiquement une expérience d’apprentissage interactive complète incluant des supports de cours, des quiz, des simulations et des activités collaboratives. Cette interaction est assurée par plusieurs agents spécialisés (enseignant, assistant et apprenants simulés) reproduisant une dynamique de classe réaliste.

Sur le plan technique, OpenMAIC repose sur une architecture d’orchestration d’agents et est compatible avec plusieurs fournisseurs de modèles de langage. En tant que solution open source, elle peut être déployée sur une infrastructure locale, offrant ainsi un contrôle total sur les données et leur confidentialité.
\subsubsection{Moodle}

Moodle \cite{moodle} est la plateforme LMS open source la plus utilisée dans le monde, 
déployée dans plus de 240 pays par des universités, des gouvernements et des entreprises. 
Elle offre un ensemble complet de fonctionnalités couvrant la création de cours, la gestion 
des évaluations, les forums de discussion, le suivi de progression et le support SCORM. 
Sa bibliothèque de plus de 1\,800 plugins permet d'étendre ses capacités avec des modules 
de gamification, d'apprentissage adaptatif et d'intégrations tierces. La version 4.5, 
publiée en octobre 2024, introduit un sous-système d'intelligence artificielle ouvrant 
la voie aux fonctionnalités de tutorat intelligent et de recommandation de contenus. 
Entièrement auto-hébergeable, Moodle offre un contrôle total sur les données et représente 
un choix privilégié pour les organisations soumises à des contraintes réglementaires strictes.

% =====================================================




\subsection{Synthèse comparative}

L'analyse des solutions présentées met en évidence des différences importantes selon les 
critères qui sont au cœur de notre projet. Les plateformes commerciales offrent des 
fonctionnalités avancées mais à des coûts élevés et avec un contrôle limité sur les données, 
tandis que les solutions open source garantissent la souveraineté des données mais manquent 
de fonctionnalités intelligentes intégrées. Aucune des solutions étudiées ne réunit à elle 
seule le coaching IA, la gamification, la génération automatique de cours, le mentorat 
structuré et une architecture multi-tenant dans un seul environnement adapté au contexte 
d'une entreprise.

\begin{table}[H]
\centering
\caption{Comparaison des solutions existantes}
\label{tab:comparaison}

\renewcommand{\arraystretch}{1.6}

\resizebox{\textwidth}{!}{%
\begin{tabular}{|p{3.8cm}|p{3cm}|p{3cm}|p{3cm}|p{3cm}|p{3cm}|p{3cm}|}

\hline

\textbf{Critère} 
& \textbf{\shortstack{Google\\Skills}} 
& \textbf{\shortstack{Duolingo}} 
& \textbf{360Learning} 
& \textbf{\shortstack{CYPHER\\Learning}}
& \textbf{OpenMAIC}
& \textbf{Moodle} \\

\hline

Coaching IA 
& Non 
& Oui (GPT-4) 
& Partiel 
& Non  
& Oui (multi-agents + LLM externe requis)
& Non \\

\hline

Gamification 
& Oui 
& Oui 
& Oui 
& Oui (pilotée par IA) 
& Non 
& Via plugin \\

\hline

Génération auto de cours 
& Non 
& Oui (IA générative) 
& Oui (à partir de PDF/DOCX) 
& Oui (CYPHER Copilot) 
& Oui (multi-agents)
& Non \\

\hline

Agents conversationnels 
& Non 
& Oui (Lily) 
& Non 
& Non 
& Oui (AI Teacher + assistants + classmates orchestrés)
& Non \\

\hline

Mentorat structuré 
& Non 
& Non 
& Partiel 
& Non  
& Non
& Non \\

\hline

Suivi des compétences 
& Partiel 
& Non 
& Oui 
& Oui 
& Non
& Partiel \\

\hline

Confidentialité des données 
& Faible (serveurs Google) 
& Faible (serveurs Duolingo) 
& Bonne (RGPD, AI Act) 
& Moyenne (cloud éditeur)  
& Totale (auto-hébergeable)
& Totale (auto-hébergeable) \\

\hline

Architecture multi-tenant 
& Non 
& Non 
& Non 
& Oui (plan Enterprise) 
& Non
& Partiel \\

\hline

Open source 
& Non 
& Non 
& Non 
& Non 
& Oui (AGPL-3.0, self-hosted, extensible, BYO LLM requis)
& Oui (GPL) \\

\hline

\textbf{Coût} 

& \textbf{Gratuit (Starter) ; Pro à 29\$/mois ; Career Certificates à 49\$/mois \cite{googleSkills}}

& \textbf{Gratuit ; Super Duolingo à partir de 12,99\$/mois ; Max avec fonctionnalités IA avancées \cite{duolingoMax}}

& \textbf{8\$/utilisateur/ mois (Team) ; Business et Enterprise sur devis \cite{learning360prix}}

& \textbf{Sur devis (à partir d’environ 10\,000\$/an selon les besoins de l’organisation) \cite{cypherLearningprix}}  

& \textbf{Gratuit (open source) ; nécessite API LLM externe et infrastructure d’hébergement \cite{openmaic}}

& \textbf{Gratuit (open source) ; coûts liés à l’hébergement et au support selon le déploiement \cite{moodle}} \\

\hline

\end{tabular}%
}

\end{table}







 
\subsection{Améliorations et innovations}

L'analyse des solutions existantes montre clairement qu'aucune d'entre elles ne répond 
de manière complète aux besoins d'une entreprise souhaitant centraliser la formation, 
le suivi des compétences et le mentorat dans un seul outil intelligent. Cette section 
présente les principaux axes d'amélioration qui guident la conception de notre plateforme.

\textbf{Mentorat structuré et coaching intelligent.}
Aucune des plateformes étudiées ne propose un accompagnement humain et intelligent 
combiné dans un même outil. Les collaborateurs en phase d'intégration se retrouvent 
souvent sans suivi structuré, tandis que les rares outils de recommandation existants 
manquent de mémoire contextuelle et ne s'adaptent pas au profil évolutif de l'apprenant. 
Notre solution place le coaching IA au centre de l'expérience d'apprentissage : un agent 
intelligent suit en continu le parcours de chaque utilisateur, identifie ses lacunes et 
lui propose un accompagnement personnalisé à tout moment. Le mentorat vient renforcer 
cette dynamique en s'appuyant sur les données du coaching pour orienter les échanges 
entre mentor et apprenant, assurer un suivi cohérent de la progression et adapter 
les recommandations aux besoins réels de chaque collaborateur.

\textbf{Génération de cours adaptée au contexte interne.}
Les outils de création automatique de cours existants produisent des contenus à partir de documents génériques. Notre solution va plus loin en permettant de générer des parcours d'apprentissage directement à partir des ressources internes de l'entreprise, en tenant compte des compétences cibles, du niveau de l'apprenant et des besoins métiers spécifiques de chaque filiale du Groupe 3S.

\textbf{Gamification au service de la motivation.}
Notre plateforme intègre des mécanismes de gamification — points, badges, niveaux et tableaux de progression — directement liés aux compétences acquises, afin de rendre l'apprentissage plus engageant et d'encourager les apprenants à aller au bout de leur parcours.

\textbf{Architecture multi-tenant adaptée au Groupe 3S.}
Parmi les solutions étudiées, seul CYPHER Learning propose une architecture multi-tenant, 
et uniquement dans son plan Enterprise à coût élevé. Notre plateforme est conçue nativement 
avec cette architecture, permettant à chaque filiale du Groupe 3S de disposer de son propre 
espace cloisonné, avec des configurations, des contenus et des utilisateurs indépendants, 
le tout sur une infrastructure partagée et mutualisée.

\textbf{Confidentialité des données et contrôle des contenus.}
Les plateformes commerciales qui intègrent des fonctionnalités d'IA (coaching, génération 
de cours, recommandations) s'appuient sur des modèles de langage externes hébergés par 
des fournisseurs tiers, ce qui implique que les contenus de formation et les données des 
utilisateurs transitent en dehors de l'organisation. Notre solution adopte une approche 
différente en utilisant des modèles de langage hébergés directement sur les serveurs de 
la société, garantissant ainsi que les contenus internes, les données des apprenants et 
les interactions avec les agents IA restent strictement confinés au sein de l'infrastructure 
de 3S, sans aucune divulgation vers l'extérieur.

\textbf{Vision unifiée du développement professionnel.}
Les outils du marché traitent séparément la formation, le suivi des compétences et le 
mentorat, obligeant les responsables à jongler entre plusieurs plateformes sans vision 
globale. Notre solution centralise ces trois dimensions dans un seul environnement cohérent, 
offrant aux managers et aux équipes RH un tableau de bord complet pour suivre l'évolution 
professionnelle de chaque collaborateur.


\section{Spécification des exigences}

Cette section présente l'ensemble des exigences identifiées pour la plateforme 
\textit{Skill Forge}. Ces exigences ont été définies à partir de l'analyse des 
besoins du Groupe 3S et des limites des solutions existantes étudiées dans la 
section précédente. Elles se divisent en deux catégories : les exigences 
fonctionnelles, qui décrivent ce que le système doit faire, et les exigences 
non fonctionnelles, qui définissent comment il doit le faire.

% =====================================================

\subsection{Exigences fonctionnelles}

Les exigences fonctionnelles décrivent les fonctionnalités que le système doit 
offrir pour répondre aux besoins des utilisateurs. Le tableau~\ref{tab:ef} 
présente l'ensemble des exigences fonctionnelles identifiées.

\renewcommand{\arraystretch}{1.3}

\begin{longtable}{|p{1.2cm}|p{4cm}|p{9.5cm}|}

\caption{Exigences fonctionnelles}
\label{tab:ef} \\

\hline
\textbf{ID} & \textbf{Catégorie} & \textbf{Description} \\
\hline
\endfirsthead

\hline
\textbf{ID} & \textbf{Catégorie} & \textbf{Description} \\
\hline
\endhead

EF-01 & Authentification & Le système doit permettre l'authentification sécurisée des utilisateurs. \\
\hline

EF-02 & Gestion des rôles & Le système doit gérer plusieurs rôles : super administrateur, administrateur, manager, expert pédagogique, responsable d’équipe et membre d’équipe. \\
\hline

EF-03 & Gestion des profils & Le système doit permettre la création et la modification des profils utilisateurs. \\
\hline

EF-04 & Multi-tenant & Le système doit permettre la gestion de plusieurs entreprises (tenants) de façon cloisonnée. \\
\hline

EF-05 & Gestion des formations & Le système doit permettre la création, la modification, la suppression et la consultation des formations. \\
\hline

EF-06 & Compétences & Le système doit permettre la gestion des compétences associées aux formations et aux utilisateurs. \\
\hline

EF-07 & Parcours d'apprentissage & Le système doit permettre la création de parcours d'apprentissage personnalisés et l'affectation de compétences cibles. \\
\hline

EF-08 & Suivi de progression & Le système doit permettre le suivi de l'avancement de chaque apprenant dans son parcours. \\
\hline

EF-09 & Recommandation IA & Le système doit recommander des contenus de formation adaptés au profil et aux besoins de chaque utilisateur. \\
\hline

EF-10 & Coaching IA & Le système doit intégrer un agent IA capable d'assister l'utilisateur tout au long de son apprentissage. \\
\hline

EF-11 & Génération de cours & Le système doit permettre la génération automatique de cours à partir de documents ou de sujets fournis. \\
\hline

EF-12 & Évaluations & Le système doit permettre la création et la gestion de quiz, tests et évaluations de validation. \\
\hline

EF-13 & Gamification & Le système doit permettre l'attribution de points, badges et récompenses selon la progression des apprenants. \\
\hline

EF-14 & Mentorat & Le système doit permettre l'affectation de mentors aux apprenants et la gestion de leurs interactions. \\
\hline

EF-15 & Tableau de bord & Le système doit fournir des tableaux de bord de suivi de la progression, des compétences et de l'engagement. \\
\hline

EF-16 & Recherche & Le système doit permettre la recherche et la consultation des formations disponibles. \\
\hline

EF-17 & Gestion des contenus & Le système doit permettre la gestion des ressources pédagogiques (vidéos, documents, liens). \\
\hline

EF-18 & Traçabilité & Le système doit enregistrer les activités des utilisateurs sur la plateforme. \\
\hline

\end{longtable}


\subsection{Exigences non fonctionnelles}

Les exigences non fonctionnelles définissent les critères de qualité que le système 
doit respecter, indépendamment des fonctionnalités qu'il offre. Conformément à la 
norme ISO/IEC 25010 \cite{iso25010}, ces exigences couvrent des dimensions telles 
que la sécurité, la performance, la maintenabilité et l'ergonomie. Le 
tableau~\ref{tab:enf} présente l'ensemble des exigences non fonctionnelles retenues.

\renewcommand{\arraystretch}{1.3}

\begin{longtable}{|p{1.6cm}|p{3cm}|p{9.5cm}|}

\caption{Exigences non fonctionnelles}
\label{tab:enf} \\

\hline
\textbf{ID} & \textbf{Catégorie} & \textbf{Description} \\
\hline
\endfirsthead

\hline
\textbf{ID} & \textbf{Catégorie} & \textbf{Description} \\
\hline
\endhead

ENF-01 & Sécurité & La plateforme doit protéger les données des utilisateurs et garantir un contrôle d'accès strict selon les rôles. \\
\hline

ENF-02 & Confidentialité & La plateforme doit assurer l'isolation complète des données entre les différents tenants. Les modèles de langage utilisés doivent être hébergés sur les serveurs de la société afin de garantir qu'aucun contenu interne ne soit transmis à des services externes. \\
\hline

ENF-03 & Performance & La plateforme doit offrir des temps de réponse acceptables et supporter plusieurs utilisateurs simultanés sans dégradation notable des performances. \\
\hline

ENF-04 & Disponibilité & La plateforme doit assurer une continuité de service fiable pour l'ensemble des utilisateurs. \\
\hline

ENF-05 & Évolutivité & La plateforme doit être conçue pour permettre l'ajout de nouvelles fonctionnalités sans remettre en cause l'architecture existante. \\
\hline

ENF-06 & Modularité & La plateforme doit reposer sur une architecture modulaire facilitant la maintenance et l'intégration de nouveaux services. \\
\hline

ENF-07 & Maintenabilité & Le code doit être clair, documenté et organisé de façon à faciliter les interventions futures. \\
\hline

ENF-08 & Extensibilité & La plateforme doit être conçue pour s'adapter à de nouveaux besoins dans un contexte multi-tenant en évolution. \\
\hline

ENF-09 & Ergonomie & La plateforme doit offrir une interface simple, claire et intuitive, accessible à des utilisateurs de profils variés. \\
\hline

ENF-10 & Accessibilité & La plateforme doit être accessible depuis différents types d'écrans et navigateurs modernes. \\
\hline

ENF-11 & Intégrité des données & La plateforme doit garantir la cohérence et l'intégrité des données à tout moment. \\
\hline

ENF-12 & Journalisation & La plateforme doit permettre le suivi des erreurs et la journalisation des événements système. \\
\hline

ENF-13 & Compatibilité IA & La plateforme doit être compatible avec l'intégration de modèles de langage et d'agents IA hébergés localement. \\
\hline

ENF-14 & Déployabilité & La plateforme doit faciliter son déploiement et la mise à jour de ses composants de manière indépendante. \\
\hline

\end{longtable}



\section{Conclusion}

Dans ce chapitre, nous avons présenté le contexte général du projet ainsi que l’entreprise d’accueil 3S et ses principaux domaines d’activité. 
Nous avons également introduit l’étude de l’existant ainsi que les différentes exigences du projet. 
Le chapitre suivant sera consacré à l’analyse et à la conception de la solution proposée.
